\documentclass{article}
\usepackage{amsmath}
\usepackage{bm}  % 替换目标宏包
\usepackage{silence}
\WarningFilter{latexfont}{Font shape}		% 屏蔽所有以 "Font shape" 开头的 LaTeX 字体警告
\WarningFilter{latexfont}{Some font shapes}	% 屏蔽所有以 "Some font shapes" 开头的 LaTeX 字体警告

\title{Test Document for \texttt{\textbackslash bf} to \texttt{\textbackslash bm} Conversion}
\author{Test Author}

\begin{document}

\maketitle

\section{Basic Usage Patterns}

% 场景1：标准的 LaTeX 2.09 样式（应被替换）
This is a {\bf bold text} in the middle of a sentence.

% 场景2：数学模式中的向量（应被替换为\bm）
Let ${\bf x}$ be a vector and ${\bf y}$ be another vector.
The sum ${\bf x} + {\bf y} = {\bf z}$.

% 场景3：带空格的变化形式
Here is {\bf another} example and {\bf  yet another  } with extra spaces.

\section{Edge Cases and Traps}

% 场景4：注释中的 \bf（应忽略，或根据你的需求决定）
% This is a comment with {\bf bold} text inside.

% 场景5：Verbatim 环境（绝对不应替换）
\begin{verbatim}
	{\bf This is verbatim text}
	\bf should remain as is
\end{verbatim}

% 场景6：行内代码（不应替换）
Use the command \verb|{\bf text}| to make bold text.

% 场景7：类似但不同的命令（不应被替换）
This uses modern \textbf{bold} and \textbf{another} approach.
Also available: {\bfseries bold series} and \textbf{with braces}.

\section{Complex Mathematical Expressions}

% 场景8：矩阵和向量（应替换）
$$
{\bf A} = \begin{pmatrix}
	{\bf a}_{11} & {\bf a}_{12} \\
	{\bf a}_{21} & {\bf a}_{22}
\end{pmatrix}
$$

% 场景9：下标中的使用（应替换）
The vector ${\bf x}_i$ has components $x_{{\bf i}j}$.

% 场景10：多重嵌套（需谨慎处理）
${\bf \hat{x}}$ and ${\bf \tilde{y}}$ with accents.

\section{Commands and Definitions}

% 场景11：宏定义中的使用（应替换）
\newcommand{\myvec}{{\bf v}}
\newcommand{\mybold}[1]{{\bf #1}}

% 场景12：段落开始
{\bf This paragraph starts with bold.} And continues normally.

% 场景13：表格中的使用
\begin{tabular}{ll}
	\hline
	{\bf Column 1} & {\bf Column 2} \\
	\hline
	{\bf Item} & Value \\
	\hline
\end{tabular}

% 多行匹配测试
Here is a very long line that contains {\bf bold text 
spanning multiple lines} which is tricky to handle.

% 转义字符测试
The sequence \textbackslash{\bf is not a command} but looks like one.

\end{document}
